
\subsection{Intro}
Key problems to solve
\begin{enumerate}
\item Chattering
\item Control gain uncertainty and overestimation
\item Input gain uncertainty 
\end{enumerate}


\subsubsection{System}
Given a single-input uncertain nonlinear system on the form
\begin{equation}
    \dot{x} = f(x,t) + h(x,t)u
    \label{eq:single-input}
\end{equation}
where $x \in \R^n$ is the state vector, $u \in \R$ is a control function, $f(x,t) \in \R$ is a differentiable, partially know vector-field.
    % \begin{block}{Note}
    %   Please note the  $|\phi(x,t)| \leq \delta|\sigma|^{1/2}$
    % \end{block}

\subsubsection{Assumptions}
\begin{assumption}
A sliding variable $\sigma = \sigma(x,t) \in \R$ is designed so that the system's \eqref{eq:single-input} desirable compensated dynamics are achieved in the sliding mode $\sigma = \sigma(x,t) = 0$.
\end{assumption}

\begin{assumption}
    The system's \eqref{eq:single-input} input-output ($u \rightarrow \sigma$) dynamics are of a relative degree one, and the internal dynamics are stable.
        Therefore, the input-output dynamics can be presented
        \begin{equation}
          \begin{split}
            \dot{\sigma} &= \underbrace{\frac{\partial\sigma}{\partial t} +  \frac{\partial\sigma}{\partial x}f(x,t)}_{\phi(x,t)} + \underbrace{\frac{\partial\sigma}{\partial x}h(x)u}_{b(x,t)} \\
            \dot{\sigma} &= \phi(x,t) + \omega,\quad \omega = b(x,t)u \iff u = b^{-1}(x,t)\omega
          \end{split}
          \label{eq:sliding-dynamics}
        \end{equation}
        The solution of system \eqref{eq:sliding-dynamics} is understood in the sense of Filippov \cite{Fil:88}.
\end{assumption}
\begin{rmk}
    We write down 
    \begin{equation*}
        \dot{\sigma} = \phi(x,t) + b(x,t)u
    \end{equation*}
\end{rmk}
\begin{rmk}
    Sense of Filippov, from Sliding Mode Control and Observation \cite{shtessel2014sliding}
\end{rmk}
For a discontinuous, or an otherwise continuous system with discontinuous control, we can write the differential equation as:
\begin{equation*}
    \dot{x} = f^c(x)
\end{equation*}
where $f^c$ is discontinuous with respect to $x$. 

\begin{wrapfigure}{r}{0.3\textwidth}
    \vspace{-50pt}
    \begin{center}
        \includegraphics[width=0.28\textwidth]{figures/filippov.png}
    \end{center}
    \vspace{-40pt}
\end{wrapfigure} 

Main idea behind the solutions proposed by Filippov for differential equations with discontinuous right-hand sides is to construct a solution as the \say{average} of the solutions obtained by approaching the point of discontinuity from different angles.

The average solution is chosen to be tangent to the manifold $\mathcal{S}$ at the point of discontinuity.
\begin{assumption}
    Further, assume that function $b(x,t) \in \R$ is known and $b(x,t) \neq 0 \ \forall x \text{ and } t \in [0, \infty)$. 
\end{assumption}
\begin{assumption}
    the perturbation function $\phi(x, t) \in \R$ is bounded 
        \begin{equation}
          |\phi(x,t)| \leq \delta |\sigma|^{1/2}
          \label{eq:shtessel-a4}
        \end{equation} 
        where the finite boundary $\delta > 0$ exists but is not known.
\end{assumption}
\begin{rmk}
    Means the perturbations are assumed to be vanishing with the square root of the sliding variable.
\end{rmk}
    % \pause
    % \begin{block}{Note}
    %   Having the constant $\delta$ as a finite boundary on the sliding mode dynamics is a big assumption
    % \end{block}

\subsection{The adaptive STW controller}
  \subsubsection{Controller}
       The STW controller as written by Shtessel et al. \cite{Sht:10} is given by
       \begin{equation}
          \begin{split}
            \omega &= -\alpha|\sigma|^{1/2}sgn(\sigma) + v \\
            \dot{v} &= -\beta sgn (\sigma)
          \end{split}
          \label{eq:shtessel-stw}
       \end{equation}
       where $\alpha = \alpha(\sigma, \dot{\sigma}, t)$ and $\beta = \beta(\sigma, \dot{\sigma}, t)$ are adaptive gains. The closed-loop sliding dynamics are then on the form
       \begin{equation}
        \begin{split}
          \dot{\sigma} &= -\alpha|\sigma|^{1/2}sgn(\sigma) + v + \phi(x,t) \\
          \dot{v} &= -\beta sgn(\sigma)
          \label{eq:shtessel-closed-loop}
        \end{split}
       \end{equation}
\begin{rmk}
    Write down closed-loop dynamics
\end{rmk}

\subsubsection{Adaptive gains}
    \begin{theorem}[Shtessel et al. \cite{Sht:10}]
      \label{th:sta}
      Consider system \eqref{eq:shtessel-closed-loop}. Suppose that the perturbation $\phi(x,t)$ satisfies \eqref{eq:shtessel-a4} for some unknown constant $\delta > 0$. Then for any initial conditions $x(0), \sigma(0)$ the sliding surface $\sigma = 0$ will be reached in finite time via STW control \eqref{eq:shtessel-stw} with the adaptive gains  
      \begin{equation}
        \begin{split}
          \dot{\alpha} &= 
          \begin{cases}
            &\omega_1 \sqrt{\frac{\gamma_1}{2}}, \text{ if } \sigma \neq 0, \\
            &0, \qquad\ \ \text{ if } \sigma = 0
          \end{cases} \\
          \beta &= 2\epsilon\alpha + \lambda + 4\epsilon^2
        \end{split}
        \label{eq:adaptation-law}
      \end{equation}
      where $\epsilon, \lambda, \gamma_1 \text{ and } \omega_1$ are arbitrary positive constants.
    \end{theorem}

\begin{rmk}
    I don't have a good answer to how they came up with this, but we can go through the proof together and see the justification for it
\end{rmk}

\subsubsection{Proof of Theorem \ref{th:sta}}
      Proof in 2 steps
      \begin{enumerate}
        \item Change of variable
        \item Lyapunov stability analysis to prove finite time convergence to the sliding manifold
        % \item<3-> Prove boundedness of adaptive gains
        % \item<4-> Prove finite 
      \end{enumerate}
      
\begin{proof}%[Theorem \ref{th:sta}]
To simplify the Lyapunov analysis we perform a change of variable
    \begin{equation*}
        z = \begin{bmatrix}
        z_1 \\ z_2 
        \end{bmatrix}
        = \begin{bmatrix}
        |\sigma|^{1/2}sgn(\sigma) \\ v
        \end{bmatrix}
    \end{equation*}
    then rewrite \eqref{eq:shtessel-closed-loop} to 
    \begin{equation}
        \begin{cases}
        \dot{z}_1 &= \frac{1}{|z_1|}\left(-\frac{\alpha}{2}z_1 + \frac{1}{2}z_2 + \frac{1}{2}\phi(x,t)\right) \\
        \dot{z}_2 &= -\frac{\beta}{|z_1|}z_1
        \end{cases}
        \label{eq:shtessel-transformed-system}
    \end{equation}

    Equation \eqref{eq:shtessel-transformed-system} can then be written on matrix form
        \begin{equation}
          \dot{z} = 
          % A(z_1)z+g(z_1)\phi(x,t)
          % \begin{bmatrix}
          %   \dot{z}_1 \\ \dot{z}_2 
          % \end{bmatrix}
          % = 
          \underbrace{
            \frac{1}{|z_1|}
          \begin{bmatrix}
            -\frac{\alpha}{2} && \frac{1}{2} \\ 
            -\beta && 0
          \end{bmatrix}
          }_{A(z_1)}
          \begin{bmatrix}
            z_1 \\ z_2 
          \end{bmatrix}
          +
          \underbrace{
            \frac{1}{|z_1|}
          \begin{bmatrix}
            \frac{1}{2} \\ 0
          \end{bmatrix} 
          }_{g(z_1)}
          \phi(x, t)
        \end{equation}

        % where 
        % \begin{equation*}
        %   z = \begin{bmatrix}
        %     \dot{z}_1 \\ \dot{z}_2 
        %   \end{bmatrix}, \
        %   A(z_1) = 
        %   \frac{1}{|z_1|}\begin{bmatrix}
        %     -\frac{\alpha}{2} && \frac{1}{2} \\ 
        %     -\beta && 0
        %   \end{bmatrix}, \
        %   g(z_1) =
        %   \frac{1}{|z_1|}\begin{bmatrix}
        %     \frac{1}{2} \\ 0
        %   \end{bmatrix}
        % \end{equation*}

        \begin{rmk}[Observe]
        We observe the following
        \begin{enumerate}
            \item If $z_1, z_2 \to 0$ in finite time then $\sigma_1, \sigma_2 \to 0$ in finite time, and
            \item $|z_1|=|\sigma|^{1/2}$ and $sgn(z_1) = sgn(\sigma)$
        \end{enumerate}
        \end{rmk}

        \textbf{Lyapunov analysis:}
        In the second step of the proof, the stability analysis of the system given by the above equation is performed.
        We introduce the LFC
        \begin{equation}
          V(z, \alpha, \beta) = V_0 + \frac{1}{2\gamma_1}\left(\alpha - \alpha^*\right)^2 + \frac{1}{2\gamma_2}\left(\beta - \beta^*\right)^2
        \end{equation}
        where 
        \begin{equation}
          V_0(z) = z^\top P z = \left(\lambda+4\epsilon^2\right)z_1^2 + z_2^2 - 4\epsilon z_1 z_2
        \end{equation}
        \begin{equation*}
          P = 
          \begin{bmatrix}
            \lambda+4\epsilon^2 && -2\epsilon \\ 
            -2\epsilon && 1
          \end{bmatrix},
          \ \forall \lambda, \epsilon > 0.
        \end{equation*}
        where $\alpha^*, \beta^* > 0$ are some constants.

        The derivative of the Lyapunov function candidate is
       \begin{equation}
        \dot{V}(z, \alpha, \beta) = \dot{z}^\top P z + z^\top P \dot{z} + \frac{1}{\gamma_1}\left(\alpha - \alpha^*\right)\dot{\alpha} + \frac{1}{\gamma_2} \left(\beta - \beta^*\right)\dot{\beta}
       \end{equation}
       We can bound the first two terms
       \begin{equation}
        \dot{V}_0 = \dot{z}^\top P z + z^\top P \dot{z} \leq -\frac{1}{|z_1|}z^\top\tilde{Q}z
       \end{equation}
       where the symmetric matrix $\tilde{Q}$ is given by
       \begin{equation}
        \tilde{Q} - 2\epsilon I = 
        \begin{bmatrix}
          2\lambda\alpha + 4\epsilon\left(2\epsilon\alpha - \beta \right) - 2\left(\lambda + 4\epsilon^2\right)\delta && * \\
          \left(\beta - 2\epsilon\alpha - \lambda - 4\epsilon^2\right) + 2\epsilon\delta && 4\epsilon 
        \end{bmatrix}
       \end{equation}
       \begin{rmk}
        (The term $- 2\epsilon I$ is likely there to simplify the right-hand side)
       \end{rmk}
       To ensure positive definite $\tilde{Q}$ we enforce
        \begin{equation}
          \beta = 2\epsilon\alpha + \lambda + 4\epsilon^2
        \end{equation}
        The matrix $\tilde{Q}$ is positive definite with a minimal eigenvalue $\lambda_{\text{min}}(\tilde{Q}) \geq 2\epsilon$ if 
        \begin{equation}
          \alpha > \frac{\epsilon\delta^2 + \left(\lambda + 4\epsilon^2\right)\left(2\epsilon + \delta\right) + \epsilon}{\lambda}
          \label{eq:shtessel-inequality-19}
        \end{equation}
        \begin{rmk}
            There would be a lot of calculations to confirm this, so I haven't bothered
        \end{rmk}
        We then have that
        \begin{equation}
          \dot{V}_0 \leq - \frac{1}{|z_1|}z^\top\tilde{Q}z \leq - \frac{\epsilon}{|z_1|}z^\top z = - \frac{\epsilon}{|z_1|}\norm{z}^2
        \end{equation}
        \begin{rmk}
            I don't quite understand where ther $- \frac{\epsilon}{|z_1|}z^\top z$ term comes from
        \end{rmk}
        We also have that  
        \begin{equation}
          \lambda_{\text{min}}\left(P\right)\norm{z}^2 \leq \underbrace{z^\top P z}_{V_0(z)} \leq \lambda_{\text{max}}\left(P\right)\norm{z}^2 
        \end{equation}
        where $\norm{z}^2 = z_1^2 + z_2^2 = |\sigma| + z_2^2$ and
        \begin{equation}
          |z_1| = |\sigma|^{1/2} \leq \norm{z}^2 \leq \frac{V_0^{1/2}(z)}{\lambda^{1/2}_{\text{min}}(P)}
        \end{equation}
        \begin{rmk}
            $\lambda^{1/2}_{\text{min}}(P)\norm{z}^2 \leq V_0^{1/2}(z)$, is this true? Shouldn't it be $\lambda^{1/2}_{\text{min}}(P)\norm{z} \leq V_0^{1/2}(z)$
        \end{rmk}
        then 
        \begin{equation}
          \dot{V}_0(z) \leq -rV_0^{1/2}, \ r=\frac{2\epsilon\lambda^{1/2}_{\text{min}}(P)}{\lambda_{\text{max}}(P)}
        \end{equation}
        \begin{rmk}
            Jeg skjønner ikke helt hvor $r$ kommer fra heller
        \end{rmk}
        We also have that  
        \begin{equation}
          \begin{split}
            \dot{V}(z, \alpha, \beta) &= \dot{z}^\top P z + z^\top P \dot{z} + \frac{1}{\gamma_1}\left(\alpha - \alpha^*\right)\dot{\alpha} + \frac{1}{\gamma_2} \left(\beta - \beta^*\right)\dot{\beta} \\
            &\leq -\frac{1}{|z_1|}z^\top\tilde{Q}z + \frac{1}{\gamma_1}\left(\alpha - \alpha^*\right)\dot{\alpha} + \frac{1}{\gamma_2} \left(\beta - \beta^*\right)\dot{\beta} \\
            &\leq -rV_0^{1/2} + \frac{1}{\gamma_1}\left(\alpha - \alpha^*\right)\dot{\alpha} + \frac{1}{\gamma_2} \left(\beta - \beta^*\right)\dot{\beta} \\
            &= -rV_0^{1/2} - \frac{\omega_1}{\sqrt{2\gamma_1}}\left|\alpha - \alpha^*\right|\dot{\alpha} - \frac{\omega_2}{\sqrt{2\gamma_2}} \left|\beta - \beta^*\right| \\
            &+ \frac{1}{\gamma_1}\left(\alpha - \alpha^*\right)\dot{\alpha} + \frac{1}{\gamma_2} \left(\beta - \beta^*\right)\dot{\beta} + \frac{\omega_1}{\sqrt{2\gamma_1}}\left|\alpha - \alpha^*\right| + \frac{\omega_2}{\sqrt{2\gamma_2}} \left|\beta - \beta^*\right|
          \end{split}
        \end{equation}
        Using the inequality
        \begin{equation}
          \left(x^2 + y^2 + z^2\right)^{1/2} \leq |x| + |y| + |z|
        \end{equation}
        using this for the Lyapunov candidate function
        \begin{equation}
          -rV_0^{1/2} - \frac{\omega_1}{\sqrt{2\gamma_1}}\left|\alpha - \alpha^*\right|\dot{\alpha} - \frac{\omega_2}{\sqrt{2\gamma_2}} \left|\beta - \beta^*\right| \leq -\eta V(z,\alpha,\beta)^{1/2}
          \label{eq:shtessel-bound-a}
        \end{equation}
        where $\eta = \text{min}(r, \omega_1, \omega_2)$.

        Using equation \eqref{eq:shtessel-bound-a} we can rewrite the upper bound on $\dot{V}$  
        \begin{equation}
          \begin{split}
            \dot{V}(z, \alpha, \beta) \leq -\eta V(z,\alpha,\beta)^{1/2} &+ \frac{1}{\gamma_1}\left(\alpha - \alpha^*\right)\dot{\alpha} + \frac{1}{\gamma_2} \left(\beta - \beta^*\right)\dot{\beta} \\
            &+ \frac{\omega_1}{\sqrt{2\gamma_1}}\left|\alpha - \alpha^*\right| + \frac{\omega_2}{\sqrt{2\gamma_2}} \left|\beta - \beta^*\right|
          \end{split}
        \end{equation}
        Assuming the adaptation law \eqref{eq:adaptation-law} makes the adaptive gains $\alpha(t)$ and $\beta(t)$ bounded there exists positive constants $\alpha^*$ and $\beta^*$ such that $\alpha(t) - \alpha^* < 0$ and $\beta(t) - \beta^* < 0, \ \forall t \geq 0$.
        Using this assumption we can reduce the upper bound expression to
        \begin{equation}
            \dot{V}(z, \alpha, \beta) \leq -\eta V(z,\alpha,\beta)^{1/2} + \xi
            \label{eq:shtessel-eq-29}
        \end{equation}
        where 
        \begin{equation}
          \xi = -\left|\alpha - \alpha^*\right|\left(\frac{1}{\gamma_1}\dot{\alpha} - \frac{\omega_1}{\sqrt{2\gamma_1}}\right) - \left|\beta - \beta^*\right|\left(\frac{1}{\gamma_2}\dot{\beta} - \frac{\omega_2}{\sqrt{2\gamma_2}}\right) 
        \end{equation}
        We want to ensure $\xi = 0$
        \begin{rmk}
            $\dot{\alpha} = \omega_1\frac{\gamma_1}{\sqrt{2\gamma_1}} = \omega_1\frac{\gamma_1\cdot\sqrt{\gamma_1}}{\sqrt{2\gamma_1}\cdot\sqrt{\gamma_1}}
            = \omega_1\frac{\gamma_1\cdot\sqrt{\gamma_1}}{\sqrt{2}\gamma_1}
            = \omega_1\sqrt{\frac{\gamma_1}{2}} $ 
        \end{rmk}
        \begin{equation}
          \dot{\alpha} = \omega_1 \sqrt{\frac{\gamma_1}{2}}
          \label{eq:shtessel-alpha-dot}
        \end{equation}
        \begin{equation}
          \dot{\beta} = \omega_2 \sqrt{\frac{\gamma_2}{2}}
        \end{equation}
        If we choose $\epsilon = \frac{\omega_2}{2\omega_1} \sqrt{\frac{\gamma_2}{\gamma_1}}$ and differentiating $\beta$ as given by \eqref{eq:adaptation-law} we get 
        \begin{equation}
          \dot{\beta} = 2\epsilon\dot{\alpha} = 2\epsilon\omega_1 \sqrt{\frac{\gamma_1}{2}} = \omega_2 \sqrt{\frac{\gamma_2}{2}}
        \end{equation}

        \textbf{Observe.}
        For the controller to converge in finite time the inequality \eqref{eq:shtessel-inequality-19} must hold. This implies that $\alpha(t)$ must increase according to \eqref{eq:shtessel-alpha-dot} until \eqref{eq:shtessel-inequality-19} is met in order to guarantee a positive definite $\tilde{Q}$ matrix.

        \begin{quote}
          Also, as soon as the sliding variable $\sigma$ and its derivative converges to zero it does not make sense to increase $\alpha(t)$ and $\beta(t)$ by making $\dot{\alpha} = 0$ as $\sigma = 0$. - Shtessel et al. \cite{Sht:10}
        \end{quote}
        Hence, we obtain the adaptation law \eqref{eq:adaptation-law}. $\qed$
        % \begin{alertblock}{Comment}
    %   It might not make sense to \emph{increase} gains after the sliding manifold is reached, however, for tracking problems it might be desirable to \emph{decrease} the control gains after the sliding manifold is reached to avoid overestimation. 
    % \end{alertblock}
\end{proof}

\subsection{Propositions}
\begin{proposition}
    The adaptive gains $\alpha(t)$ and $\beta(t)$ are bounded.
\end{proposition}
\begin{proof}
    A solution to the adaptive gain $\alpha(t)$ in \eqref{eq:adaptation-law}
    \begin{equation}
    \alpha =
    \begin{cases}
        \alpha(0) + \omega_1\sqrt{\frac{\gamma_1}{2}}t, \ \ 0 \leq t \leq t_r \\
        \alpha(0) + \omega_1\sqrt{\frac{\gamma_1}{2}}t_r, \ t > t_r
    \end{cases}
    \end{equation}
    Hence, $\alpha(t)$ and $\beta(t)$ are bounded. $\qed$
\end{proof}

\begin{proposition}
    The finite time $t_r$ can be estimated as
\begin{equation}
    t_r \leq \frac{2V^{1/2}(0)}{\eta}
\end{equation}
where $\eta = \text{min}(r, \omega_1, \omega_2)$.
\end{proposition}
\begin{proof}
    Since the right-hand side of the inequality \eqref{eq:shtessel-inequality-19} is bounded and $\alpha(t)$ is increasing linearly with respect to time according to the adaptation law \eqref{eq:adaptation-law}, it is met in finite time.
      The above inequality is found by integrating \eqref{eq:shtessel-eq-29}.
      $\qed$
\end{proof}

Key takeaways
\begin{itemize}
    \item Deals well with matched uncertainties
    \item Chattering is attenuated by having the switching term behind an integrator
    \item Adaptive gain law is obtained and drives the system to the desired state in finite time
\end{itemize}
\begin{itemize}
    \item The input gain is assumed known
    \item Controller gains are increasing
    \item Many design variables with tuning lack of intuition
\end{itemize}
